\SourcePage{606}{609}
\section{General theory of scattering}

In classical mechanics, a collision between two particles is completely specified by their velocities and impact parameter, the distance by which their paths would miss one another in the absence of interaction. Quantum mechanics changes the formulation of the question itself: for motion with specified velocities, a trajectory has no meaning, and consequently neither does an impact parameter. The theory seeks only the probability that the collision deflects the particles---or, as one says, \emph{scatters} them---through a given angle. Here we consider so-called elastic collisions, in which the particles undergo no transformations and, if they are composite, their internal states do not change.

Like every two-body problem, an elastic collision reduces to the scattering of one particle of reduced mass by the field $U(r)$ of a fixed force centre\footnote{We neglect the spin--orbit interaction of the particles, if they possess spin. By taking the field to be centrally symmetric, we also exclude processes such as the scattering of electrons by molecules.}. The reduction is made by passing to the coordinate frame in which the centre of mass of the two particles is at rest. Denote the scattering angle in this frame by $\theta$. Simple formulas relate it to the deflection angles $\theta_1$ and $\theta_2$ in the ``laboratory'' frame, where one of the particles, the second, was initially at rest:
\begin{equation}
 \tan\theta_1=\frac{m_2\sin\theta}{m_1+m_2\cos\theta},\qquad
 \theta_2=\frac{\pi-\theta}{2},
 \tag{123.1}
\end{equation}
where $m_1$ and $m_2$ are the particle masses (see I, \S~17). In particular, for equal masses, $m_1=m_2$, this becomes simply
\begin{equation}
 \theta_1=\frac\theta2,\qquad \theta_2=\frac{\pi-\theta}{2};
 \tag{123.2}
\end{equation}
\SourcePage{607}{610}
thus $\theta_1+\theta_2=\pi/2$, and the particles depart at right angles.

Throughout the remainder of this chapter, unless expressly stated otherwise, we use the centre-of-mass frame and take $m$ to mean the reduced mass of the colliding particles.

A free particle travelling in the positive $z$ direction is described by a plane wave, written $\psi=e^{ikz}$. This selects a normalization in which the current density in the wave equals the particle speed $v$. Far from the centre, the scattered particles are described by a diverging spherical wave $f(\theta)e^{ikr}/r$, where $f(\theta)$ is a function of the scattering angle $\theta$, the angle between the $z$ axis and the scattered-particle direction. This function is called the \emph{scattering amplitude}. The exact wave function solving the Schrödinger equation with potential energy $U(r)$ must therefore have the large-distance asymptotic form
\begin{equation}
 \psi\approx e^{ikz}+\frac{f(\theta)}r e^{ikr}.
 \tag{123.3}
\end{equation}

The probability per unit time for a scattered particle to pass through the surface element $dS=r^2do$, where $do$ is an element of solid angle, is $vr^{-2}|f|^2dS=v|f|^2do$\footnote{This argument tacitly assumes that the incident particle beam is bounded by a broad aperture, to avoid diffraction effects, but one of finite size, as in actual scattering experiments. There is consequently no interference between the two terms in (123.3): $|\psi|^2$ is evaluated at points where the incident wave is absent.}. Its ratio to the current density of the incident wave is
\begin{equation}
 d\sigma=|f(\theta)|^2do.
 \tag{123.4}
\end{equation}
This quantity has the dimensions of area and is called the \emph{effective cross-section}, or simply the \emph{cross-section}, for scattering into the solid-angle element $do$. Setting $do=2\pi\sin\theta\,d\theta$ gives the cross-section
\begin{equation}
 d\sigma=2\pi\sin\theta|f(\theta)|^2d\theta
 \tag{123.5}
\end{equation}
for scattering into the angular interval from $\theta$ to $\theta+d\theta$.

The solution of the Schrödinger equation describing scattering in a central field $U(r)$ must clearly be axially symmetric about the $z$ axis, the incident direction. Any such solution can be written as a\SourcePage{608}{611} superposition of continuous-spectrum wave functions for particles moving in the field with fixed energy $\hbar^2k^2/2m$, with orbital angular momenta of various magnitudes $l$ and zero $z$ projections. These functions do not depend on the azimuthal angle $\varphi$ about the $z$ axis and are therefore axially symmetric. The required wave function consequently has the form
\begin{equation}
 \psi=\sum_{l=0}^\infty A_lP_l(\cos\theta)R_{kl}(r),
 \tag{123.6}
\end{equation}
where the $A_l$ are constants and the radial functions $R_{kl}(r)$ obey
\begin{equation}
 \frac1{r^2}\frac d{dr}\left(r^2\frac{dR_{kl}}{dr}\right)
 +\left[k^2-\frac{l(l+1)}{r^2}-\frac{2m}{\hbar^2}U(r)\right]R_{kl}=0.
 \tag{123.7}
\end{equation}
The coefficients $A_l$ must be chosen so that (123.6) has the asymptotic form (123.3) at large distances. We shall show that the required choice is
\begin{equation}
 A_l=\frac1{2k}(2l+1)i^l\exp(i\delta_l),
 \tag{123.8}
\end{equation}
where $\delta_l$ are the phase shifts of the functions $R_{kl}$. This also solves the problem of expressing the scattering amplitude in terms of those phases.

The asymptotic form of $R_{kl}$ is given by (33.20):
\[
 \begin{split}
 R_{kl}&\approx\frac2r\sin\left(kr-\frac{l\pi}{2}+\delta_l\right)\\
 &=\frac1{ir}\{(-i)^l\exp[i(kr+\delta_l)]
 -i^l\exp[-i(kr+\delta_l)]\}.
 \end{split}
\]
Substitution of this expression and (123.8) into (123.6) gives the asymptotic wave function
\begin{equation}
 \psi\approx\frac1{2irk}\sum_{l=0}^\infty(2l+1)P_l(\cos\theta)
 [(-1)^{l+1}e^{-ikr}+S_le^{ikr}],
 \tag{123.9}
\end{equation}
where we have introduced the notation
\begin{equation}
 S_l=\exp(2i\delta_l).
 \tag{123.10}
\end{equation}
On the other hand, expansion (34.2) of the plane wave, after the same transformation, is
\[
 e^{ikz}\approx\frac1{2ikr}\sum_{l=0}^\infty(2l+1)P_l(\cos\theta)
 [(-1)^{l+1}e^{-ikr}+e^{ikr}].
\]

\SourcePage{609}{612}
In the difference $\psi-e^{ikz}$ all terms containing $e^{-ikr}$ cancel, as they must. The coefficient of $e^{ikr}/r$ in that difference, namely the scattering amplitude, is therefore
\begin{equation}
 f(\theta)=\frac1{2ik}\sum_{l=0}^\infty(2l+1)(S_l-1)P_l(\cos\theta).
 \tag{123.11}
\end{equation}
This expresses the scattering amplitude through the phases $\delta_l$ (H.~Faxen, J.~Holtsmark, 1927)\footnote{The inverse problem of reconstructing the scattering potential from phases $\delta_l$ assumed known is of fundamental interest. It was solved by \emph{I.~M.~Gelfand, B.~M.~Levitan, and V.~A.~Marchenko}. In principle, determining $U(r)$ requires knowledge of $\delta_0(k)$ as a function of the wave vector over the entire range from $k=0$ to $k=\infty$, together with the coefficients $a_n$ in the asymptotic expressions
\[
 R_{n0}\approx a_ne^{-\kappa_nr}/r
 \quad(\kappa_n=\sqrt{2m|E_n|}/\hbar)
\]
as $r\to\infty$ for the wave functions of any states belonging to discrete, negative energy levels $E_n$. From these data, $U(r)$ is found by solving a particular linear integral equation. A systematic account is given in V.~de~Alfaro and T.~Regge, \emph{Potential Scattering}, Moscow: Mir, 1966.}.

Integration of $d\sigma$ over all angles gives the total scattering cross-section $\sigma$, the ratio of the total probability per unit time that a particle is scattered to the current density in the incident wave. Substituting (123.11) in
\[
 \sigma=2\pi\int_0^\pi|f(\theta)|^2\sin\theta\,d\theta
\]
and using the mutual orthogonality of Legendre polynomials with distinct $l$, together with
\[
 \int_0^\pi P_l^2(\cos\theta)\sin\theta\,d\theta=\frac2{2l+1},
\]
we obtain the total cross-section
\begin{equation}
 \sigma=\frac{4\pi}{k^2}\sum_{l=0}^\infty(2l+1)\sin^2\delta_l.
 \tag{123.12}
\end{equation}

Each term of this sum is a \emph{partial cross-section} $\sigma_l$ for scattering of particles with a specified orbital\SourcePage{610}{613} angular momentum $l$. Its greatest possible value is
\begin{equation}
 \sigma_{l\,\max}=\frac{4\pi}{k^2}(2l+1).
 \tag{123.13}
\end{equation}
Comparison with (34.5) shows that the number of particles scattered with angular momentum $l$ can be four times the number of such particles in the incident flux. This is a purely quantum effect caused by interference between scattered and unscattered particles.

It will also be convenient below to use the \emph{partial scattering amplitudes} $f_l$, defined as the coefficients in the expansion
\begin{equation}
 f(\theta)=\sum_{l=0}^\infty(2l+1)f_lP_l(\cos\theta).
 \tag{123.14}
\end{equation}
By (123.11), they are related to the phases $\delta_l$ by
\begin{equation}
 f_l=\frac1{2ik}(S_l-1)=\frac1{2ik}(e^{2i\delta_l}-1),
 \tag{123.15}
\end{equation}
and the partial cross-sections are
\begin{equation}
 \sigma_l=4\pi(2l+1)|f_l|^2.
 \tag{123.16}
\end{equation}

\ProblemHeading
Express the scattering amplitude through phase shifts in two dimensions. The field is $U=U(\rho)$, where $\rho=\sqrt{x^2+z^2}$, and the particle flux is directed along the $z$ axis.

\SolutionHeading In two dimensions, the wave function far from the scatterer is a superposition of a plane wave and a diverging cylindrical wave:
\begin{equation}
 \psi=e^{ikz}+f(\varphi)\frac{e^{ik\rho}}{\sqrt{-i\rho}}.
 \tag{1}
\end{equation}
Here $\varphi$ is the angle between the $z$ axis and the direction of scattering, while $f(\varphi)$ is the scattering amplitude, whose dimension in two dimensions is the square root of length. The factor $-i=\exp(-i\pi/2)$ under the radical is introduced to simplify the formulas that follow. The scattering cross-section per unit length along the $y$ axis is
\[
 d\sigma=|f|^2d\varphi.
\]
It has the dimension of length.

The wave function must be expanded in functions with a definite projection $m$ of angular momentum on the $y$ axis, of the form $Q_m(\rho)e^{im\varphi}$. Far from the scatterer, the radial functions differ from the free-motion functions obtained in the problem to \S~34 only by a phase shift:
\[
 Q_m(\rho)\approx i^m\sqrt{\frac2{\pi k\rho}}
 \sin\left[k\rho-\frac\pi2\left(m-\frac12\right)+\delta_m\right],
\]
\SourcePage{611}{614}
where $\delta_m=\delta_{-m}$. Repeating the argument of this section and using the plane-wave expansion from the problem to \S~34, we find that the function with asymptotic form (1) is given by the series
\[
 \psi=\sum_{m=-\infty}^{\infty}e^{i\delta_m}Q_m(\rho)e^{im\varphi},
\]
and the scattering amplitude is
\begin{equation}
 f(\varphi)=\frac1{i\sqrt{2\pi k}}
 \sum_{m=-\infty}^{\infty}(e^{2i\delta_m}-1)e^{im\varphi}.
 \tag{2}
\end{equation}
Integration gives the total cross-section
\[
 \sigma=\int_0^{2\pi}|f|^2d\varphi
 =\sum_{m=-\infty}^{\infty}\sigma_m,
 \qquad \text{where}\qquad
 \sigma_m=\sigma_{-m}=\frac4k\sin^2\delta_m.
\]
One readily verifies the relation
\begin{equation}
 \operatorname{Im}f(0)=\sqrt{\frac{k}{8\pi}}\,\sigma,
 \tag{3}
\end{equation}
which is the optical theorem in two dimensions (see formula (125.9) below).
